\begin{longtable}{@{}>{\raggedright\arraybackslash}p{4.4cm}>{\raggedright\arraybackslash}p{6.5cm}>{\raggedright\arraybackslash}p{3.95cm}@{}}
\caption{Training exposure, prior transfer and generic-task retention.}\label{tab:s2}\\
\toprule
Model & Quantity & Count \\
\midrule\endfirsthead
\multicolumn{3}{l}{\small Table \thetable{} (continued)}\\\toprule
Model & Quantity & Count \\
\midrule\endhead
\bottomrule\endfoot
CellPFN & Parameters & 12329889 \\
CellPFN & Optimizer updates & 96000 \\
CellPFN & Simulated banks & 12000 \\
CellPFN & Scheduled episodes & 384000 \\
CellPFN & Loss-bearing episodes & 276828 \\
TabICLv2 & Parameters & 27552258 \\
TabICLv2 & Additional cellular tasks & 40000 \\
TabICLv2 & Adaptation updates & 2500 \\
Kumo small & Parameters & 27458266 \\
Kumo small & Additional cellular tasks & 40000 \\
Kumo small & Adaptation updates & 2500 \\
Kumo large & Parameters & 213668250 \\
Kumo large & Additional cellular tasks & 40000 \\
Kumo large & Adaptation updates & 2500 \\
\end{longtable}
CellPFN banks are reused for multiple episodes; scheduled and loss-bearing episodes are distinct. Adaptation uses 40,000 cellular tasks and 10,000 generic replay tasks per arm. Released-model pretraining precedes this adaptation and is not included in these counts. TabICLv2 reports 35 million synthetic pretraining datasets\citep{qu2026tabiclv2}; this published count is distinct from CellPFN banks and scheduled episodes. These units do not establish a matched-compute comparison.
\medskip\noindent\textbf{Adapted minus released balanced accuracy (percentage points).}\par
\nopagebreak[4]\par\noindent\begin{tabular}{@{}>{\raggedright\arraybackslash}p{4.4cm}>{\raggedright\arraybackslash}p{3.5cm}>{\raggedright\arraybackslash}p{3.5cm}>{\raggedright\arraybackslash}p{3.45cm}@{}}
\toprule
Model & 1 support & 10 supports & 100 supports \\
\midrule
TabICLv2 & +7.30 & -0.50 & +0.77 \\
Kumo small & +10.79 & +1.31 & -0.22 \\
Kumo large & +5.58 & +1.08 & +0.26 \\
\bottomrule
\end{tabular}\par
Each delta uses released weights re-scored in the same matched scoring run as the adapted checkpoint. Table~\ref{table:overview} instead reports the stored benchmark scores for released models; subtracting its displayed means need not reproduce these paired deltas.\par
\medskip\noindent\textbf{Kumo-small low-support change relative to released weights.}\par
\nopagebreak[4]\par\noindent\begin{tabular}{@{}>{\raggedright\arraybackslash}p{7.5cm}>{\raggedright\arraybackslash}p{7.35cm}@{}}
\toprule
Training stream & Mean change (percentage points) \\
\midrule
Cellular task stream 1 & +6.05 \\
Cellular task stream 2 & +5.73 \\
Cellular task stream 3 & +5.17 \\
Generic-table control & -0.05 \\
\bottomrule
\end{tabular}\par
The low-support average weights eight panel--budget combinations equally (three support draws each). The three cellular runs vary the task stream with the same model seed. The generic control changes rows and columns per update, so this is not a matched-data or matched-cost contrast. No corresponding generic-only control was run for TabICLv2 or Kumo large.\par
\medskip\noindent\textbf{Retention on four generic classification datasets.}\par
\nopagebreak[4]\par\noindent\begin{tabular}{@{}>{\raggedright\arraybackslash}p{4.1cm}>{\raggedright\arraybackslash}p{2.65cm}>{\raggedright\arraybackslash}p{2.65cm}>{\raggedright\arraybackslash}p{2.65cm}>{\raggedright\arraybackslash}p{2.8cm}@{}}
\toprule
Model & One-shot BA & One-shot accuracy & All-budget BA & All-budget accuracy \\
\midrule
TabICLv2 released & 70.9 & 69.3 & 87.8 & 87.2 \\
TabICLv2 adapted & 64.0 & 62.0 & 84.8 & 84.3 \\
\bottomrule
\end{tabular}\par
Balanced accuracy (BA) and ordinary accuracy (\%) equally weight breast cancer, digits, iris and wine, with five split seeds each. The all-budget mean additionally weights 1, 10 and the full training split equally. These exploratory retention checks do not estimate universal tabular performance.